\scanpage{43}\mspage{49}
On applique \uline{Der \deleted{5}4} (g) et \uline{Der \deleted{6}5} (g)

\deleted{\uline{Cor 2} (g) Soit \unsure{$S$|$\mathscr S$}
\diagram{\begin{tikzcd}[ampersand replacement=\&,row sep=small]
 X'\arrow[rr,"\alpha"]\arrow[dr,"f'"'] \& \& X\arrow[dl,"f"]\\
 \& S \&
\end{tikzcd}}
comm. dans Diag, avec $\alpha$ hot-équiv. stricte. Alors $\forall F$ dans $\mathbb D(X)$, on a
\[
 \alpha^*:f_*(F)\xrightarrow{\sim} f'_*(\alpha^*(F))
\]}

% slanted
\margin{Critère de changt de base pour images directes cohomologiques}
\uline{Cor. \unsure{3|5}} (g) Soit
\diagram{\begin{tikzcd}[ampersand replacement=\&,column sep=large]
 X'\arrow[r,"q"]\arrow[d,"f'"'] \& X\arrow[d,"f",""{name=right,coordinate}]\\
 Y'\arrow[r,"p",""{name=bottom,coordinate}] \& Y
 \arrow[from=right,to=bottom,bend right=12,"\unsure{\gamma|r}"']
\end{tikzcd}}
dans Diag (cf page 30), d'où hom de chgt de base
\[
 p^*f_*(F)\longrightarrow f'_*q^*F\qquad(F\in\mathrm{Ob}\,\mathbb D(X)).
\]
Pour que ce soit un iso. pour tout $F$, il suffit que \deleted{\unsure{le}} $\forall y'\in\mathrm{Ob}\,Y'$, le foncteur induit par $(p,q)$
\[
 X'_{/y'}\longrightarrow X_{/p(y')}
\]
une \uline{hot-équivalence stricte}\inserted{ (ou costricte)} ; cela signifie aussi
% slanted
\margin{\uline{Faux} pour ``stricte'', ne marche que pour ``\uline{co}stricte''}
\inserted{(\deleted{\unsure{dans}} \inserted{dans} le cas strict)} que pour
\[
 y'_0\in\mathrm{Ob}\,Y',\quad x_0\in\mathrm{Ob}\,X,\quad u_0:f(x_0)\to p(y'_0)
\]
la catégorie \inserted{$K(y'_0,x_0,u_0)$} des objets de la forme
\[
 (x',v',u)\mid x'\in\mathrm{Ob}\,X',\quad v':f'(x')\to y'_0,\quad u:q(x')\to x_0
\]
rendant commutatif le carré dans $Y'$
\diagram{$(*)$\quad\begin{tikzcd}[ampersand replacement=\&,column sep=large]
 f(q(x'))\arrow[r,"\gamma(x')"]\arrow[d,"f(u)"'] \& p(f'(x'))\arrow[d,"p(v')"]\\
 f(x_0)\arrow[r,"u_0"']\arrow[u,dashed,shift right=1ex,"\text{\scriptsize en cas de hoto-équiv. costricte}"'] \& p(y'_0)
\end{tikzcd}}
% slanted
\margin{comparer p. 30 \unsure{ff.}}
% slanted
\margin{\uline{NB} Dans le cas de la hot. équivalence \uline{costricte}, c'est la \unsure{=|même} description de $K={}_{(x_0,u_0)\backslash}(X'_{/\unsure{y'_0}})$, sauf que la flèche $u$ va en sens inv., $x_0\to q(x')$, et idem dans le diagramme \unsure{carré} $(*)$}
est \unsure{contractile|contractible}.

On applique \deleted{\unsure{Dér 5 (g)}} \inserted{Dér. 2, Der 4(g)} et \uline{Cor 1} \unsure{$\geq$|$\leq$|et}
% archive sheet number: 42
